\begin{ledgroup}
\footnotesize
\pstart
\noindent
In den nachfolgenden Stücken N.\;\ref{60458} und N.\;\ref{41200} widmet sich Leibniz zwei Dissertationen \indexN{\textls{Spoleti}, Francesco, gest. 1712}Francesco Spoletis:
\cite{02095}\textit{De momento quo gravia nituntur deorsum per datum planum} \lbrack...\rbrack\ \textit{de secretione bilis in Hepate}, Venedig 1686.
Bei N.\;\ref{60458} handelt es sich um Auszüge mit Bemerkungen, vorwiegend aus der ersten Dissertation;
bei der Reinschrift N.\;\ref{41200} (aus der Hand J.\,D.~Brandshagens\indexN{\textls{Brandshagen} (Brondhaguen), Jobst Dietrich 1659\textendash nach 1716})
lassen Form und Inhalt auf die Textgattung Rezension schließen.
Diese Besprechung, die beide Werke in ähnlichem Umfang behandelt, blieb nach heutigem Wissensstand unveröffentlicht (die anonyme \cite{02117}Rezension in den \cite{01023}\textit{Acta Eruditorum} 1687, S.\;627--629 stammt von \indexN{\textls{Knorr} (Knorre), Martin 1657\textendash1699}Martin Knorr).
Als Vorlage der Reinschrift ist zwar ein eigenhändiges Konzept anzunehmen, jedoch ist kein solches bekannt.
\pend

\pstart
Spoletis erste Dissertation, statischen Inhalts, widmet sich der
um 1684 von \indexN{\textls{Vanni}, Giovanni Francesco 1638\textendash1709}G.\,F. Vanni entfachten
und besonders in Italien heftig ausgetragenen
Kontroverse über die korrekte Zerlegung der Kräfte auf schiefen Ebenen (wozu auch Leibniz publizierte, siehe N.\;\ref{55606} und N.\;\ref{49409}).
Der Dissertation war eine Flugschrift vorausgegangen:
\indexN{\textls{Spoleti}, Francesco, gest. 1712}\textsc{Spoleti}, \cite{02096}\textit{De momento gravis in plano inclinato}, Rom 1684.
\pend

\pstart
Den Terminus post quem für die Entstehung der Auszüge (damit auch der Rezension)
bietet die Erscheinung des Werks um 1686.
Der vorgeschlagene Terminus ante quem \glqq Mitte Februar 1687\grqq\ 
ergibt sich aus dem Briefwechsel mit \indexN{\textls{Mencke} (Menken, Menkenius, Menque), Otto 1644\textendash1707}O.\;Mencke
und aus dem Verhältnis der Reinschrift N.\;\ref{41200} mit der zweiten auf demselben Bogen überlieferten Reinschrift, d.\,h. dem Textzeugen \textit{l} von N.\;\ref{41199} (Leibnizens Rezension zu
\indexN{\textls{Vanni}, Giovanni Francesco 1638\textendash1709}\textsc{G.\,F. Vanni}, \cite{02000}\textit{Exegeses physico-mathematicae}).
Der Schreibduktus und die Textverteilung auf dem Träger lassen auf eine gemeinsame, konsekutive Abfassung beider Texte durch J.\,D.~Brandshagen\indexN{\textls{Brandshagen} (Brondhaguen), Jobst Dietrich 1659\textendash nach 1716} schließen: erst N.\;\ref{41199}, dann N.\;\ref{41200}.
Die Reinschrift des Textes über Vanni samt Leibnizens Verbesserungen (Textschicht \textit{Lil}  von N.\;\ref{41199})
bildete die Vorlage für die in den \cite{01023}\textit{Acta} von April 1687
gedruckte Rezension (Textzeuge \textit{E}; siehe die Datierung von N.\;\ref{41199}).
Da die Eintragung der Verbesserungen nach der (als etwa gleichzeitig angenommene) Abfassung beider Reinschriften erfolgt sein muss,
darf der Terminus ante quem von N.\;\ref{41199} (Zusendung der Rezension an \indexN{\textls{Mencke} (Menken, Menkenius, Menque), Otto 1644\textendash1707}Mencke am 21. Februar 1687) auf N.\;\ref{41200} übertragen werden.
\pend
\pstart
Der Bogen LH~III 5 Bl.~17\textendash18 überliefert neben N.\;\ref{60458} ein weiteres Stück (Überschrift "De secretione animali"). Es handelt sich dabei um umfangreiche kommentierte Auszüge aus \indexN{\textls{Spoleti}, Francesco, gest. 1712}Spoletis Dissertation \textit{De secretione bilis in hepate}, die vsl. in einem späteren Band der Reihe erscheinen werden.
Der Text ist wie folgt auf dem Träger verteilt: Anfang von "De secretione animali" auf Bl.~17~r\textsuperscript{o}, Fortsetzung auf Bl.~17~v\textsuperscript{o}, Schlussteil in dicht gedrängter Schrift ebd. am unteren, dann linken Rand; auf Bl.~18~r\textsuperscript{o} unser Stück N.\;\ref{60458}; Bl.~18~v\textsuperscript{o} leer.
Dieser Umstand legt nahe, dass Leibniz zunächst auf Bl.~18~r\textsuperscript{o} die Auszüge N.\;\ref{60458}
anfertigte, die die Dissertation \textit{De secretione} nur beiläufig behandeln (lediglich wenige Zeilen zu den ersten drei Seiten der Vorlage);
als er dann, auf Grund seines großen Interesses an der Dissertation, auf Bl.~17~r\textsuperscript{o} die Exzerpierung erneut aufnahm (ab der zweiten Seite des Werks und viel ausführlicher), musste er in zunehmend kleiner und dicht gedrängter Schrift fortfahren, da Bl.~18~r\textsuperscript{o} nun schon beschrieben war.
\pend
\pstart
Leibnizens Interesse an \indexN{\textls{Spoleti}, Francesco, gest. 1712}Spoletis innovativem Forschungsansatz belegt neben seiner Bemerkung im Schlussatz von N.\;\ref{41200} (\refpassage{LH_35_09_20_003-004_2a}{LH_35_09_20_003-004_2b})
auch der Bericht über sein persönliches Gespräch mit \indexN{\textls{Spoleti}, Francesco, gest. 1712}Spoleti 
in \indexO{Venedig}Venedig
im \cite{02135}Brief an \indexN{\textls{Bianchini}, Francesco 1672\textendash1729}Bianchini vom 18. März 1690:
\glqq A \indexN{\textls{Spoleti}, Francesco, gest. 1712}Spoleto non contemnenda expecto, juvenis enim est, et ingenio valens, et praeclarae quantum judicare possum voluntatis. Hortatus sum, ut mathematicum in re medica agat, quoad ejus fieri potest, id enim hactenus mea sententia desideratur, et demto \indexN{\textls{Borelli} (Borellus), Giovanni Alfonso 1608\textendash1679}Borello, qui medicus non fuit, magis observationes aut divinationes quam firmas ratiocinationes circa corpora animalium scriptores nobis dedere\grqq\ 
(\cite{02135}\textit{LSB} III,~4 N.\;244, S.\;481 Z.~13--17).
\pend
\end{ledgroup}
